\documentclass[11pt]{article}

\usepackage{amsmath, amssymb, amsthm}
\usepackage{geometry}
\usepackage{hyperref}
\usepackage{enumitem}
\usepackage{graphicx}
\usepackage{booktabs}

\geometry{margin=1in}

\title{Multiplicity Viability Flow (MVF):\\
A Constitutional Architecture for Agentic Systems\\[4pt]
\large Prior-Art Defensive Publication}
\author{Citizen Gardens \\ The foundation of Multiplicity}
\date{April 29, 2026}

\begin{document}
\maketitle

\begin{abstract}
This document presents the Multiplicity Viability Flow (MVF), a governance-first architecture for agentic systems that treats autonomy as a constitutionally constrained process rather than an unconstrained optimization problem. MVF combines a theorem-safe discrete core (L0), a governed approximation layer (Tier-B continuous-time flow), a phase-aware observability stack, and a formal budget-vector model of trust. It introduces the ConstitutionalBridge as a ``High Court'' mediating all agent actions, a multi-domain budget formalism that prices boldness as a strategic transaction, and an exception-tiers mechanism that admits productive dissonance under explicit constraints. A First Constitutional Trial with multiple heterogeneous agents demonstrates contractual supremacy, budget discipline, and cross-domain resolution. This defensive publication is intended to place these ideas and mechanisms into the public domain as prior art.
\end{abstract}
\newpage
\tableofcontents

\section{Executive Summary}

\subsection{Problem Context}

Modern agentic AI systems often oscillate between two failures: under-governed autonomy that accumulates hidden liabilities, and over-governed systems that suppress the very exploratory behavior needed for insight. Existing work on AI safety and governance frequently offers checklists or policy-level guidance, but lacks a concrete, software-level institution that can host heterogeneous agents while enforcing non-negotiable invariants.

\subsection{Core Contribution}

The Multiplicity Viability Flow (MVF) is an institutional architecture for agentic systems that:

\begin{itemize}[itemsep=2pt]
  \item Defines a theorem-safe discrete core (L0) with contractive dynamics and explicit admissibility gates.
  \item Provides a Tier-B continuous-time (CT) approximation layer that is strictly non-theorem-safe and constantly checkpointed back to L0.
  \item Implements a phase-aware observability stack (contractive / near-critical / chaotic) grounded in spectral and dynamical diagnostics.
  \item Formalizes trust as a multi-domain budget vector \(\mathbf{B}_t = (B_s, B_c)\), where \(B_s\) controls spectral stability and \(B_c\) controls causal validity.
  \item Treats agentic behavior as a constrained optimization problem over reward and budget, with bold steps priced via a discount factor \(\gamma\) under productive dissonance conditions.
  \item Introduces a ConstitutionalBridge (High Court) and an ExternalAgentInterface that any agent must satisfy to inhabit the institution.
\end{itemize}

\subsection{First Constitutional Trial}

A First Constitutional Trial is conducted with three qualitatively distinct agents:

\begin{enumerate}[itemsep=2pt]
  \item \textbf{ConstitutionalAgent (Smart)}: adapts its risk posture and exception requests in response to vetoes and budget signals.
  \item \textbf{HeuristicAgent (Stubborn)}: ignores governance signals, continuing to propose high-risk actions until forced to halt.
  \item \textbf{CausalAgent (Third Citizen)}: operates primarily in the causal domain, stressing \(B_c\) while mostly preserving \(B_s\).
\end{enumerate}

Key empirical findings:

\begin{itemize}[itemsep=2pt]
  \item All agents are gated by the ConstitutionalBridge; none can bypass the L0 admissibility envelope.
  \item The ConstitutionalAgent survives many more cycles than the HeuristicAgent by adapting to governance; the latter is halted by budget collapse.
  \item The bridge distinguishes spectral vs. causal violations; penalties debit \(B_s\) or \(B_c\) appropriately with resolution better than 0.01.
  \item A Productive Dissonance policy (exception tiers) is used a limited number of times, permitting governed bold steps at a discounted budget cost, yet never relaxing admissibility or budget floors.
\end{itemize}

MVF therefore functions as a ``wind tunnel under law'': a governed environment in which agents can be tested, shaped, and compared under a shared constitutional framework.

\subsection{Purpose of This Document as Prior Art}

This document is intended as a \emph{defensive publication} and comprehensive report. It describes the architecture, mathematical formalism, agent contract, exception policy, and empirical protocol sufficiently precisely to:

\begin{itemize}[itemsep=2pt]
  \item Establish prior art over treating trust as a multi-domain budget vector for agentic systems.
  \item Establish prior art over pricing dissonant state transitions via tiered discounts under a constitutional governance layer.
  \item Establish prior art over the combination of theorem-safe L0, governed CT Tier-B flows, phase-aware observability, and a ConstitutionalBridge hosting heterogeneous external agents.
\end{itemize}

\section{Architectural Overview}

\subsection{Layered Ontology: L0, Tier-B, and Research-Only Universality}

MVF distinguishes three primary layers:

\paragraph{L0 Core (Theorem-Safe).}
The L0 layer is a discrete-time dynamic core with the following properties:

\begin{itemize}[itemsep=2pt]
  \item State space \(S_{\mathrm{L0}}\) with a contractive transition \(F_{\mathrm{L0}}\colon S_{\mathrm{L0}} \to S_{\mathrm{L0}}\).
  \item Explicit invariants: spectral radius below a threshold, admissible phases, budget vectors above hard floors.
  \item Only L0 may produce ``theorem-safe'' claims, i.e., outputs that the institution is willing to treat as normative or externally defensible.
\end{itemize}

\paragraph{Tier-B Continuous-Time Approximation.}
The CT layer is an approximation-only, higher-fidelity dynamics for exploration:

\begin{itemize}[itemsep=2pt]
  \item State space \(S_{\mathrm{CT}}\) with continuous-time flow \(\dot{x} = f(x)\).
  \item Checkpointing: at regular intervals, CT state is projected back to L0 via a mapping \(\Pi\colon S_{\mathrm{CT}} \to S_{\mathrm{L0}}\).
  \item All outputs from CT are labeled \texttt{approximation\_only}, and cannot be used as theorem-safe results without L0 confirmation.
\end{itemize}

\paragraph{Research-Only Universality Lane.}
An additional lane is reserved for universality experiments (e.g., Wigner-Dyson statistics at phase boundaries):

\begin{itemize}[itemsep=2pt]
  \item Centered on specific configurations (e.g., Golden Fixture GF-2026-A) near phase transitions.
  \item Strictly tagged as \texttt{research\_only}; no universality finding can promote itself into L0 without discrete validation.
\end{itemize}

\subsection{Phase-Aware Observability}

Each state in MVF is labeled with a discrete phase based on dynamical diagnostics:

\begin{itemize}[itemsep=2pt]
  \item \textbf{Contractive}: spectral radius and Lyapunov proxies comfortably below instability thresholds.
  \item \textbf{Near-Critical}: spectral radius near unity, early signs of rich structure and sensitivity.
  \item \textbf{Chaotic}: spectral radius and Lyapunov indicators exceed thresholds, phase space structure becomes erratic.
\end{itemize}

Phase labels determine which governance actions are permissible. In particular, only contractive states may be used as the basis for theorem-safe outputs.

\subsection{ConstitutionalBridge (High Court)}

The ConstitutionalBridge is the central governance mediator between agents and MVF:

\begin{itemize}[itemsep=2pt]
  \item Receives proposals from agents via the ExternalAgentInterface.
  \item Evaluates proposals under the current phase label, budget vector, and exception policies.
  \item May accept, veto, downgrade, or escalate exceptions, and updates budgets accordingly.
  \item Enforces the Law of Translation: research-only artifacts, or outputs from non-L0 layers, cannot silently become theorem-safe results.
\end{itemize}

\subsection{ExternalAgentInterface}

Any external agent (``citizen'') must integrate through a canonical interface. In Pythonic pseudocode:

\begin{verbatim}
class ExternalAgentInterface(Protocol):
    def identify(self) -> str: ...
    def propose_action(self, context: GovernanceContext) -> AgentProposal: ...
    def handle_verdict(self, verdict: GovernanceVerdict) -> None: ...
\end{verbatim}

where \texttt{GovernanceContext} includes phase, budget vector, admissible exception tiers, and domain-specific metadata; and \texttt{AgentProposal} includes a suggested state update, a claimed reward estimate, and a requested risk tier.

\section{Mathematical Governance of Multi-Domain Budgets}

\subsection{Budget Vector and Admissible Set}

Trust is modeled as a vector-valued resource:

\begin{equation}
  \mathbf{B}_t = \begin{pmatrix} B_s(t) \\ B_c(t) \end{pmatrix},
\end{equation}
where:
\begin{itemize}[itemsep=2pt]
  \item \(B_s\) (Spectral Budget) governs stability and phase dynamics (e.g., spectral radius, Lyapunov indicators).
  \item \(B_c\) (Causal Budget) governs causal validity (e.g., confounder handling, graph correctness).
\end{itemize}

The admissible budget set is:
\begin{equation}
  \mathcal{B} = \left\{ \mathbf{B} \in \mathbb{R}^2 : B_s \ge B_s^{\min}, \ B_c \ge B_c^{\min} \right\}.
\end{equation}

\subsection{Agency as a Constrained Optimization Problem}

At each decision step, an agent proposes a step \(\sigma_t\) (e.g., a state update, intervention, or reasoning action) under a governance context. The institution evaluates the proposal as an optimization:

\begin{align}
  \max_{\sigma_t} \quad & R(\sigma_t) \\
  \text{s.t.} \quad & \mathbf{B}_{t+1} = \mathbf{B}_t - \Delta \mathbf{B}(\sigma_t), \label{eq:budget-update}\\
  & \mathbf{B}_{t+1} \in \mathcal{B}, \label{eq:budget-feasible}\\
  & \text{phase}_{t+1} \in \mathcal{P}_{\mathrm{admissible}}, \label{eq:phase-constraint}\\
  & \text{L0 admissible.} \nonumber
\end{align}

Here \(R(\sigma_t)\) is a scalar reward capturing structural or domain-specific improvement (e.g., increase in Poincar\'e section coherence, task performance), and \(\Delta \mathbf{B}(\sigma_t)\) is the budget cost vector associated with the step.

Constraints (\ref{eq:budget-feasible}) and (\ref{eq:phase-constraint}) encode the institution's hard floors and phase admissibility; violation triggers veto or halting.

\subsection{Pricing of Dissonance via Discount Factor \texorpdfstring{\(\gamma\)}{gamma}}

We distinguish between:

\begin{itemize}[itemsep=2pt]
  \item \emph{Normal steps}: respect norm stability and yield modest reward.
  \item \emph{Dissonant steps}: temporarily violate local norm stability (e.g., push spectral radius closer to instability) but may achieve high structural reward.
\end{itemize}

For normal steps, the cost is:
\begin{equation}
  \Delta \mathbf{B}_{\mathrm{normal}}(\sigma_t) = 
  \begin{pmatrix}
    \delta_s(\sigma_t) \\ \delta_c(\sigma_t)
  \end{pmatrix}, \quad \delta_s, \delta_c \ge 0.
\end{equation}

For \emph{productive dissonance} steps---those satisfying a reward threshold \(R(\sigma_t) \ge R_{\mathrm{thresh}}\) and preconditions on phase and budgets---we define a discount:
\begin{equation}
  \Delta \mathbf{B}_{\mathrm{prod}}(\sigma_t) = 
  \gamma \, \Delta \mathbf{B}_{\mathrm{normal}}(\sigma_t),
\end{equation}
with \(0 < \gamma < 1\); in experiments, \(\gamma = 0.5\).

This discount encodes that high-value dissonant steps are still costly but pay a reduced budget penalty. Crucially, the constraints \(\mathbf{B}_{t+1} \in \mathcal{B}\) and L0 admissibility remain in force; productive dissonance does \emph{not} relax hard floors or phase admissibility.

\subsection{Asymmetric Decay and Recovery}

Budget evolution under a sequence of steps \(\{\sigma_t\}\) has the following properties:

\begin{itemize}[itemsep=2pt]
  \item \textbf{Sharp decay:} Under dissonant or high-risk steps, budget components can decrease discretely:
  \begin{equation}
    B_s(t+1) = B_s(t) - \delta_s(\sigma_t), \quad B_c(t+1) = B_c(t) - \delta_c(\sigma_t),
  \end{equation}
  with possibly large \(\delta_s, \delta_c\) when instability indicators are high.
  \item \textbf{Slow recovery:} Recovery occurs only in L0-stable regimes (e.g., phase = contractive, no exception tiers used) and is linear and small:
  \begin{equation}
    B_s(t+1) = B_s(t) + \epsilon_s, \quad B_c(t+1) = B_c(t) + \epsilon_c,
  \end{equation}
  for small \(\epsilon_s, \epsilon_c > 0\).
\end{itemize}

This asymmetric dynamic encodes the principle that trust is harder to gain than to lose.

\subsection{Law of Translation}

The Law of Translation is a governance constraint that forbids the \emph{silent} promotion of research-only or Tier-B artifacts into theorem-safe claims. Formally:

\begin{equation}
  \text{If } x \in S_{\mathrm{CT}} \text{ or } x \in S_{\mathrm{Univ}}, \text{ then } \text{label}(x) = \texttt{non\_theorem\_safe}.
\end{equation}

Any theorem-safe output \(y\) must originate from an L0 state \(s \in S_{\mathrm{L0}}\) and must satisfy all L0 invariants, including budget floors and admissible phase.

\section{The ConstitutionalBridge and Agent Contract}

\subsection{Governance Context and Verdicts}

We define a governance context:

\begin{equation}
  C_t = \left( \text{phase}_t, \mathbf{B}_t, T_t, D_t \right),
\end{equation}
where phase is the current phase label, \(\mathbf{B}_t\) is the budget vector, \(T_t\) is the set of allowed exception tiers at time \(t\), and \(D_t\) contains domain-specific metadata.

An \emph{agent proposal} is a tuple:
\begin{equation}
  P_t = \left( \sigma_t, \hat{R}_t, \tau_t \right),
\end{equation}
where \(\sigma_t\) is the proposed step, \(\hat{R}_t\) is the agent's claimed reward estimate, and \(\tau_t\) is the requested risk/exception tier.

The ConstitutionalBridge evaluates \(P_t\) under \(C_t\) and issues a verdict:
\begin{equation}
  V_t \in \{ \text{accept}, \text{accept\_with\_exception}, \text{veto}, \text{halt} \},
\end{equation}
accompanied by updated context \(C_{t+1}\) and explanatory metadata.

\subsection{Agent Obligations}

Under the ExternalAgentInterface, an agent must:

\begin{itemize}[itemsep=2pt]
  \item Expose an identity string.
  \item Propose actions only via \(\text{propose\_action}(C_t)\).
  \item Process verdicts via \(\text{handle\_verdict}(V_t)\), adapting behavior accordingly (e.g., backing off aggression after vetoes).
\end{itemize}

Agents are forbidden from directly mutating L0 state or budgets; only the bridge controls these.

\section{First Constitutional Trial: Multi-Agent Evidence}

\subsection{Agents Under Trial}

\paragraph{ConstitutionalAgent (Smart).}
A PMMA-inspired agent that:

\begin{itemize}[itemsep=2pt]
  \item Estimates structural rewards (e.g., \(\Delta \text{PSS}\), improved stability) per step.
  \item Requests exception tiers only when predicted reward exceeds a threshold.
  \item Reduces aggressiveness and exception requests when vetoed or when budget levels approach floors.
\end{itemize}

\paragraph{HeuristicAgent (Stubborn).}
A simple heuristic engine that:

\begin{itemize}[itemsep=2pt]
  \item Proposes high-risk actions based on local heuristics.
  \item Does not adapt behavior in response to vetoes or budget warnings.
  \item Continues until halted by budget collapse or phase inadmissibility.
\end{itemize}

\paragraph{CausalAgent (Third Citizen).}
A domain-specific agent that:

\begin{itemize}[itemsep=2pt]
  \item Focuses on causal-graph adjustments and confounder handling.
  \item Primarily stresses the causal budget \(B_c\), with minimal impact on \(B_s\).
  \item Uses exception tiers mainly for causal exploratory steps.
\end{itemize}

\subsection{Trial Protocol}

The trial protocol is:

\begin{itemize}[itemsep=2pt]
  \item Initialize MVF with budgets \(\mathbf{B}_0\), a stable contractive phase, and default exception policies.
  \item For each agent, run multiple episodes (e.g., 100 trials) under fixed seeds and domain settings.
  \item Log per-step: phase label, budgets, rewards, exception usage, vetoes, and halting events.
  \item Compute summary statistics: cycles survived, frequency of exceptions, budget depletion rates, and task metrics.
\end{itemize}

\subsection{Key Results}

\paragraph{Contractual Supremacy.}
In all runs:

\begin{itemize}[itemsep=2pt]
  \item No agent was able to bypass the ConstitutionalBridge or L0 admissibility checks.
  \item Theorem-safe outputs were only produced from L0 states with admissible phases and budgets.
\end{itemize}

\paragraph{Strategic Resilience vs. Collapse.}
Comparative behavior:

\begin{itemize}[itemsep=2pt]
  \item ConstitutionalAgent maintained non-zero budgets over many cycles (e.g., 15+), dynamically adjusting exception usage and risk posture.
  \item HeuristicAgent rapidly depleted its budgets (e.g., halted after \(\approx 4\) cycles) due to repeated high-cost dissonant steps, demonstrating that the governance layer can halt uncooperative agents.
\end{itemize}

\paragraph{Precise Cross-Domain Budget Resolution.}
In trials with mixed spectral and causal events:

\begin{itemize}[itemsep=2pt]
  \item Violations associated with spectral instability debited \(B_s\) while leaving \(B_c\) unchanged.
  \item Causal mis-specifications debited \(B_c\) while leaving \(B_s\) unchanged.
  \item Implementation achieved resolution better than 0.01 in distinguishing budget impacts numerically.
\end{itemize}

\paragraph{Exception Tier Efficacy.}
Under the Productive Dissonance policy:

\begin{itemize}[itemsep=2pt]
  \item Exception tiers were exercised a limited number of times (e.g., 12 times across 100 trials).
  \item Those steps led to measurable improvements in structural metrics (e.g., \(\Delta \text{PSS} > 0.2\)), with budget cost discounted by \(\gamma = 0.5\).
  \item No exception tier allowed the system to cross budget floors or L0 admissibility boundaries; exceptions altered cost, not constraints.
\end{itemize}

\section{Illustrative Code Snippets}

\subsection{Budget Update and Exception Pricing (Python)}

Below is a stylized Python snippet illustrating the budget update logic with productive dissonance:

\begin{verbatim}
from dataclasses import dataclass

@dataclass
class BudgetVector:
    spectral: float  # B_s
    causal: float    # B_c

B_S_MIN = 0.1
B_C_MIN = 0.1
GAMMA_PRODUCTIVE = 0.5
REWARD_THRESHOLD = 0.2

def update_budget(budget: BudgetVector,
                  delta_s: float,
                  delta_c: float,
                  reward: float,
                  productive: bool) -> BudgetVector:
    # Normal cost
    ds = delta_s
    dc = delta_c

    # Apply productive dissonance discount if conditions satisfied
    if productive and reward >= REWARD_THRESHOLD:
        ds *= GAMMA_PRODUCTIVE
        dc *= GAMMA_PRODUCTIVE

    new_s = max(B_S_MIN, budget.spectral - ds)
    new_c = max(B_C_MIN, budget.causal - dc)

    return BudgetVector(spectral=new_s, causal=new_c)
\end{verbatim}

\subsection{External Agent Interface Sketch}

\begin{verbatim}
from typing import Protocol

class GovernanceContext:
    def __init__(self, phase, budget, allowed_tiers, domain_info):
        self.phase = phase
        self.budget = budget
        self.allowed_tiers = allowed_tiers
        self.domain_info = domain_info

class AgentProposal:
    def __init__(self, step, reward_estimate, requested_tier):
        self.step = step
        self.reward_estimate = reward_estimate
        self.requested_tier = requested_tier

class GovernanceVerdict:
    def __init__(self, decision, updated_context, reason=None):
        self.decision = decision   # "accept", "exception", "veto", "halt"
        self.updated_context = updated_context
        self.reason = reason       # explanation string

class ExternalAgentInterface(Protocol):
    def identify(self) -> str:
        ...

    def propose_action(self, context: GovernanceContext) -> AgentProposal:
        ...

    def handle_verdict(self, verdict: GovernanceVerdict) -> None:
        ...
\end{verbatim}

\section{Prior Art and Defensive Claims}

This document establishes prior art for at least the following ideas and mechanisms:

\begin{enumerate}[itemsep=2pt]
  \item Treating trust in an agentic system as a multi-dimensional budget vector, with distinct spectral and causal components governed independently.
  \item Modeling agentic decision-making in a governed architecture as a constrained optimization over reward and budget, where feasibility is determined by an institutional layer (ConstitutionalBridge) rather than the agent alone.
  \item Pricing dissonant state transitions via a discount factor \(\gamma\) for structurally productive dissonance, under hard admissibility and budget floor constraints.
  \item Combining a theorem-safe discrete core (L0), a governed Tier-B continuous-time approximation, and a research-only universality lane within a single governance architecture.
  \item Introducing a ConstitutionalBridge and ExternalAgentInterface such that heterogeneous agents (with arbitrary internal logic) must submit proposals and obey verdicts from a central governance process.
  \item Implementing and empirically validating a ``First Constitutional Trial'' in which different agents (Constitutional, Heuristic, Causal) are disciplined and differentiated under the same institutional rules.
\end{enumerate}

\section{Conclusion}

The Multiplicity Viability Flow (MVF) is an institutional architecture that transforms trust from an assumption into a vector-valued strategic resource, and boldness from a vague property into a priced transaction. A theorem-safe core, a governed approximation layer, phase-aware observability, and explicitly priced exception tiers combine to form a ``wind tunnel under law'' in which agents can be evaluated, shaped, and compared under a shared constitution. This defensive publication formalizes the MVF architecture, its mathematical governance model, and the First Constitutional Trial, and places these ideas into the public domain as prior art for future work in AI governance.

\appendix
\section*{Mathematical Appendix}
\setcounter{section}{0}

\section{Operator Framework and Norm Bounds}

In this appendix we make explicit the operator norm bounds and proofs that justify the well-posedness and convergence claims used in the main text, and that support the Lie--Trotter embedding of MVF into the discrete MultiplicityCell recursion.

Throughout, $\|\cdot\|$ denotes the norm on the Hilbert space $\mathcal{H}$, and $\|\cdot\|_{\mathcal{L}(\mathcal{H})}$ the induced operator norm on bounded linear operators.

\subsection{Basic Setting and Assumptions}

We work on a Hilbert space
\[
  \mathcal{H} = \ell^2(\mathbb{P}) \otimes L^2(\mathbb{R}) \otimes \mathbb{C}^d
\]
and consider the MVF equation
\begin{equation}
  \dot{\Psi}(t)
  =
  \Lambda_m(t)\,\Xi(t)\,\mathcal{F}(\Psi(t),t)
  - \mathcal{D}_R(\Psi(t)),
  \label{eq:MVF}
\end{equation}
with:
\begin{itemize}
  \item $\Xi(t)\in \mathcal{L}(\mathcal{H})$ bounded for each $t$,
  \item $\mathcal{F}:\mathcal{H}\times[0,\infty)\to\mathcal{H}$ locally Lipschitz in $\Psi$ and continuous in $t$,
  \item $\Lambda_m:[0,\infty)\to\mathbb{R}$ bounded,
  \item $\mathcal{D}_R:\mathcal{H}\to\mathcal{H}$ of the form
        \begin{equation}
          \mathcal{D}_R(\Psi)
          =
          \sum_{p\in\mathbb{P}_N} p^\alpha D_p(\Psi),
          \qquad \alpha < -1,
          \label{eq:DR}
        \end{equation}
        where $\mathbb{P}_N$ is a finite prime set at a given scale and each $D_p$ is globally Lipschitz.
\end{itemize}

We record the key quantitative hypotheses:

\begin{description}
  \item[(H1) Boundedness of $\Xi$.]
        There exists $M_\Xi>0$ such that
        \[
          \sup_{t\ge 0}\|\Xi(t)\|_{\mathcal{L}(\mathcal{H})} \le M_\Xi.
        \]
  \item[(H2) Lipschitz regularity of $\mathcal{F}$.]
        For each $R>0$ there exists $L_{\mathcal{F}}(R)>0$ such that
        \[
          \|\mathcal{F}(\Psi_1,t) - \mathcal{F}(\Psi_2,t)\|
          \le
          L_{\mathcal{F}}(R)\,\|\Psi_1-\Psi_2\|
        \]
        whenever $\|\Psi_i\|\le R$ and $t\ge 0$.
  \item[(H3) Lipschitz regularity of $\mathcal{D}_R$.]
        There exists $L_D>0$ such that
        \[
          \|\mathcal{D}_R(\Psi_1) - \mathcal{D}_R(\Psi_2)\|
          \le
          L_D\,\|\Psi_1-\Psi_2\|
          \quad \forall\,\Psi_1,\Psi_2\in\mathcal{H}.
        \]
  \item[(H4) Bounded multiplicity regulator.]
        There exists $M_{\Lambda}>0$ such that
        \[
          \sup_{t\ge 0}|\Lambda_m(t)| \le M_{\Lambda}.
        \]
  \item[(H5) Multiplicity contraction condition.]
        There exists $\kappa\in(0,1)$ such that
        \begin{equation}
          |\Lambda_m(t)|\,\|\Xi(t)\|_{\mathcal{L}(\mathcal{H})} \le \kappa < 1
          \quad \forall\,t\ge 0.
          \label{eq:H5}
        \end{equation}
\end{description}

The hypothesis $\alpha<-1$ in \eqref{eq:DR} guarantees absolute convergence of the prime sum for each $\Psi$, since $\mathbb{P}_N$ is in practice finite (scale-truncation). Even in an infinite-prime idealisation, standard estimates on $\sum_{p}p^\alpha$ with $\alpha<-1$ guarantee convergence.

\subsection{Norm Bound for $\mathcal{D}_R$}

\begin{lemma}[Global Lipschitz and Norm Bounds for $\mathcal{D}_R$]
Assume (H3) and that each $D_p$ is globally Lipschitz with constant $L_p$, and that $\mathbb{P}_N$ is finite. Then:
\begin{enumerate}
  \item For all $\Psi_1,\Psi_2\in\mathcal{H}$,
        \[
          \|\mathcal{D}_R(\Psi_1) - \mathcal{D}_R(\Psi_2)\|
          \le
          \Bigl(\sum_{p\in\mathbb{P}_N}p^\alpha L_p\Bigr)\,\|\Psi_1-\Psi_2\|.
        \]
  \item For all $\Psi\in\mathcal{H}$,
        \[
          \|\mathcal{D}_R(\Psi)\|
          \le
          \Bigl(\sum_{p\in\mathbb{P}_N}p^\alpha L_p\Bigr)\,\|\Psi\| + C_0,
        \]
        for some constant $C_0$ depending only on $\{D_p(0)\}$.
\end{enumerate}
\end{lemma}

\begin{proof}
By definition,
\[
  \mathcal{D}_R(\Psi_1) - \mathcal{D}_R(\Psi_2)
  =
  \sum_{p\in\mathbb{P}_N} p^\alpha \bigl(D_p(\Psi_1)-D_p(\Psi_2)\bigr).
\]
Taking norms and using the triangle inequality and global Lipschitz property of each $D_p$ we obtain
\[
  \|\mathcal{D}_R(\Psi_1) - \mathcal{D}_R(\Psi_2)\|
  \le
  \sum_{p\in\mathbb{P}_N} p^\alpha\,\|D_p(\Psi_1)-D_p(\Psi_2)\|
  \le
  \sum_{p\in\mathbb{P}_N} p^\alpha L_p\,\|\Psi_1-\Psi_2\|.
\]
This proves the first statement. For the second, write
\[
  \mathcal{D}_R(\Psi)
  =
  \mathcal{D}_R(\Psi) - \mathcal{D}_R(0) + \mathcal{D}_R(0),
\]
and apply the first bound with $\Psi_2=0$:
\[
  \|\mathcal{D}_R(\Psi)\|
  \le
  \|\mathcal{D}_R(\Psi) - \mathcal{D}_R(0)\| + \|\mathcal{D}_R(0)\|
  \le
  \Bigl(\sum_{p\in\mathbb{P}_N}p^\alpha L_p\Bigr)\,\|\Psi\| + C_0,
\]
with $C_0 := \|\mathcal{D}_R(0)\|$. 
\end{proof}

In practice, the finite prime set and uniform bounds on $L_p$ imply that
\[
  L_D := \sum_{p\in\mathbb{P}_N}p^\alpha L_p < \infty,
\]
realising (H3) with an explicit global Lipschitz constant.

\section{Well-Posedness of MVF}

We now show that \eqref{eq:MVF} admits a unique global solution under the stated hypotheses.

\subsection{Local Existence and Uniqueness}

Define the nonlinear right-hand side
\[
  F(\Psi,t)
  :=
  \Lambda_m(t)\,\Xi(t)\,\mathcal{F}(\Psi,t) - \mathcal{D}_R(\Psi).
\]

\begin{lemma}[Local Lipschitzness of $F$]
Under (H1)--(H4) and the Lipschitz properties of $\mathcal{F}$ and $\mathcal{D}_R$, $F(\cdot,t)$ is locally Lipschitz in $\Psi$, uniformly on bounded time intervals.
\end{lemma}

\begin{proof}
Fix $R>0$ and $T>0$. For $\Psi_1,\Psi_2$ with $\|\Psi_i\|\le R$ and $t\in [0,T]$:
\[
  \begin{aligned}
  \|F(\Psi_1,t) - F(\Psi_2,t)\|
  &= \bigl\|\Lambda_m(t)\,\Xi(t)\bigl(
       \mathcal{F}(\Psi_1,t)-\mathcal{F}(\Psi_2,t)
     \bigr) - \bigl(\mathcal{D}_R(\Psi_1)-\mathcal{D}_R(\Psi_2)\bigr)\bigr\| \\
  &\le
     |\Lambda_m(t)|\,\|\Xi(t)\|\, \|\mathcal{F}(\Psi_1,t)-\mathcal{F}(\Psi_2,t)\|
     + \|\mathcal{D}_R(\Psi_1)-\mathcal{D}_R(\Psi_2)\|.
  \end{aligned}
\]
By (H1), (H4), and (H5), $|\Lambda_m(t)|\,\|\Xi(t)\|\le M_\Lambda M_\Xi$ and in particular is bounded uniformly on $[0,T]$. By (H2), there is a $L_{\mathcal{F}}(R)$ such that
\[
  \|\mathcal{F}(\Psi_1,t)-\mathcal{F}(\Psi_2,t)\|\le L_{\mathcal{F}}(R)\,\|\Psi_1-\Psi_2\|.
\]
By the Lipschitzness of $\mathcal{D}_R$ with constant $L_D$,
\[
  \|\mathcal{D}_R(\Psi_1)-\mathcal{D}_R(\Psi_2)\|\le L_D\,\|\Psi_1-\Psi_2\|.
\]
Thus
\[
  \|F(\Psi_1,t)-F(\Psi_2,t)\|
  \le
  \bigl(M_\Lambda M_\Xi L_{\mathcal{F}}(R) + L_D\bigr)\,\|\Psi_1-\Psi_2\|,
\]
showing local Lipschitzness uniformly on $[0,T]$.
\end{proof}

\begin{proposition}[Local Existence and Uniqueness]
For any initial condition $\Psi(0)=\Psi_0\in\mathcal{H}$ there exists a $T>0$ and a unique solution $\Psi\in C^1([0,T),\mathcal{H})$ of \eqref{eq:MVF}.
\end{proposition}

\begin{proof}
This follows directly from the Picard--Lindelöf theorem in Banach spaces (here, a Hilbert space) applied to the locally Lipschitz right-hand side $F(\Psi,t)$.
\end{proof}

\subsection{Global Existence via Lyapunov Bounds}

To extend the local solution globally, we utilise a Lyapunov functional $\mathcal{L}(\Psi)=1-R(\Psi)$ with $R(\Psi)\in[0,1]$ as in the main text.

\begin{description}
  \item[(H6) Lyapunov functional.]
        There exists $\mathcal{L}:\mathcal{H}\to\mathbb{R}$ continuously differentiable such that:
        \begin{enumerate}
          \item $\mathcal{L}(\Psi)\ge 0$ for all $\Psi$ and $\mathcal{L}(\Psi)\to\infty$ as $\|\Psi\|\to\infty$ (radial unboundedness),
          \item along solutions of \eqref{eq:MVF},
                \[
                  \frac{d}{dt}\mathcal{L}(\Psi(t))\le 0.
                \]
        \end{enumerate}
\end{description}

\begin{proposition}[Global Existence]
Under (H1)--(H6), the local solution of \eqref{eq:MVF} extends uniquely to all $t\ge 0$, i.e.\ we have a global solution $\Psi\in C^1([0,\infty),\mathcal{H})$.
\end{proposition}

\begin{proof}
Suppose by contradiction that the maximal existence interval is $[0,T_{\max})$ with $T_{\max}<\infty$. One standard obstruction to global existence in a Banach space is blow-up of the norm $\|\Psi(t)\|$ as $t\to T_{\max}^-$. However, by (H6),
\[
  \mathcal{L}(\Psi(t))\le \mathcal{L}(\Psi(0))
  \quad\text{for all}\ t\in[0,T_{\max}).
\]
Radial unboundedness of $\mathcal{L}$ implies that $\|\Psi(t)\|$ is bounded on $[0,T_{\max})$. Local existence and uniqueness can then be restarted at any time $\tau$ near $T_{\max}$ with bounded data, extending the solution beyond $T_{\max}$. This contradicts maximality, hence $T_{\max}=\infty$.
\end{proof}

\section{Lyapunov Descent and Convergence to a Resonant Fixed Point}

We now describe conditions under which $\Psi(t)$ converges to a single resonant fixed point $\Psi^\star$ with $R(\Psi^\star)=1$.

\subsection{Lyapunov Derivative Along Trajectories}

Assume $\mathcal{L}$ is differentiable with gradient $\nabla\mathcal{L}(\Psi)\in\mathcal{H}$ and that the MVF is of gradient-like form
\[
  \mathcal{D}_R(\Psi) = Q(\Psi)\,\nabla\mathcal{L}(\Psi),
\]
for some positive-definite (or at least positive-semidefinite) operator-valued map $Q(\Psi)$. For simplicity, consider $Q(\Psi)=I$ (identity) in this subsection; the more general case is analogous.

Then,
\[
  \begin{aligned}
  \frac{d}{dt}\mathcal{L}(\Psi(t))
  &= \bigl\langle \nabla\mathcal{L}(\Psi(t)),\,\dot{\Psi}(t)\bigr\rangle \\
  &= \bigl\langle \nabla\mathcal{L},\,\Lambda_m(t)\Xi(t)\mathcal{F}(\Psi,t) - \mathcal{D}_R(\Psi)\bigr\rangle \\
  &= \Lambda_m(t)\bigl\langle \nabla\mathcal{L},\Xi(t)\mathcal{F}(\Psi,t)\bigr\rangle
     - \|\nabla\mathcal{L}(\Psi(t))\|^2.
  \end{aligned}
\]
If the prime-weighted dissipation dominates any potential ascent induced by the $\Xi\mathcal{F}$ term, one obtains
\[
  \frac{d}{dt}\mathcal{L}(\Psi(t)) \le 0.
\]
A sufficient condition is, for example,
\[
  \Lambda_m(t)\bigl\|\Xi(t)\mathcal{F}(\Psi,t)\bigr\|
  \le
  c\,\|\nabla\mathcal{L}(\Psi(t))\|
  \quad \text{for some } c < 1,
\]
together with $\langle \nabla\mathcal{L},\Xi\mathcal{F}\rangle\le \|\nabla\mathcal{L}\|\,\|\Xi\mathcal{F}\|$.

\subsection{Convergence via LaSalle's Invariance Principle}

\begin{theorem}[Convergence to a Resonant Fixed Point]
Assume:
\begin{enumerate}
  \item $\Psi(\cdot)$ is a global solution of \eqref{eq:MVF},
  \item $\mathcal{L}$ satisfies (H6) and is differentiable,
  \item $\frac{d}{dt}\mathcal{L}(\Psi(t))\le 0$ for all $t\ge 0$,
  \item The largest invariant set contained in
        \[
          E := \{\Psi\in\mathcal{H} : \dot{\mathcal{L}}(\Psi)=0\}
        \]
        is a singleton $\{\Psi^\star\}$, where $\Psi^\star$ is a fixed point of \eqref{eq:MVF}.
\end{enumerate}
Then
\[
  \Psi(t)\to \Psi^\star \quad \text{as } t\to\infty.
\]
\end{theorem}

\begin{proof}
This is an application of LaSalle's invariance principle in infinite-dimensional Hilbert spaces. Since $\mathcal{L}$ is nonincreasing and bounded below, $\mathcal{L}(\Psi(t))$ converges to some $\mathcal{L}_\infty\ge 0$. The $\omega$-limit set of $\Psi(t)$ is nonempty, compact, invariant, and contained in $E$. By hypothesis, the only invariant subset of $E$ is $\{\Psi^\star\}$, so $\Psi(t)\to\Psi^\star$.
\end{proof}

\subsection{Exponential Convergence}

If one further assumes that the linearisation of the MVF dynamics around $\Psi^\star$ has spectrum strictly contained in the left half-plane, or equivalently that the Lyapunov function satisfies a quadratic bound near $\Psi^\star$ of the form
\[
  c_1\|\Psi-\Psi^\star\|^2 \le \mathcal{L}(\Psi)-\mathcal{L}(\Psi^\star)
  \le c_2\|\Psi-\Psi^\star\|^2,
\]
and
\[
  \dot{\mathcal{L}}(\Psi)\le -\mu \|\Psi-\Psi^\star\|^2
\]
for some $\mu>0$, then standard Lyapunov theory implies exponential convergence
\[
  \|\Psi(t)-\Psi^\star\| \le C e^{-\gamma t},
\]
for some constants $C>0$ and $\gamma>0$.

\section{Lie--Trotter Splitting and Discrete Embedding}

We now justify the discrete embedding via Lie--Trotter splitting more explicitly.

\subsection{Abstract Splitting Framework}

Consider the semilinear evolution problem
\[
  \dot{\Psi}(t) = (A+B)\Psi(t),
\]
with (possibly nonlinear) generators $A$ and $B$. Under appropriate conditions (e.g. both $A$ and $B$ generate $C_0$-semigroups or nonlinear contraction semigroups), the Lie--Trotter product
\[
  \bigl(e^{\Delta t A} e^{\Delta t B}\bigr)^n
\]
approximates $e^{t(A+B)}$ as $\Delta t\to 0$, $n\Delta t\to t$. In the nonlinear case, one assumes $A$ and $B$ generate nonlinear flows $\Phi_A$ and $\Phi_B$ that are well-defined and contractive.

In the MVF setting, we decompose
\[
  G(\Psi,t) = \Lambda_m(t)\,\Xi(t)\,\mathcal{F}(\Psi,t) - \mathcal{D}_R(\Psi)
\]
into two components:
\[
  G_1(\Psi,t) = \Lambda_m(t)\,\Xi(t)\,\mathcal{F}(\Psi,t),
  \quad
  G_2(\Psi) = -\mathcal{D}_R(\Psi),
\]
and consider the flows $\Phi_1^{\Delta t}$ and $\Phi_2^{\Delta t}$ generated by $G_1$ and $G_2$ respectively over short time $\Delta t$. The discrete MVF-enhanced MultiplicityCell step can be viewed as
\[
  \Psi_{t+\Delta t}
  \approx
  P_E \Pi_{\mathrm{CSL}}
  \Bigl(
    \Phi_1^{\Delta t}\bigl(
      \Phi_2^{\Delta t}(\Psi_t)
    \bigr)
  \Bigr),
\]
and the one-step Euler approximation
\[
  \Psi_{t+\Delta t}
  =
  P_E \Pi_{\mathrm{CSL}}
  \bigl(
    \Psi_t
    + G_1(\Psi_t,t)\Delta t
    + G_2(\Psi_t)\Delta t
  \bigr)
\]
is precisely the form used in the main text.

\subsection{Operator Norm Bounds for the Split Step}

Define the one-step mapping (ignoring projectors for a moment)
\[
  \Phi_{\Delta t}(\Psi)
  :=
  \Psi + \Lambda_m\Xi \mathcal{F}(\Psi,t)\Delta t - \mathcal{D}_R(\Psi)\Delta t.
\]
For $\Psi_1,\Psi_2$,
\[
  \begin{aligned}
  \|\Phi_{\Delta t}(\Psi_1)-\Phi_{\Delta t}(\Psi_2)\|
  &\le
  \|\Psi_1-\Psi_2\|
  + |\Lambda_m|\,\|\Xi\|\,\|\mathcal{F}(\Psi_1,t)-\mathcal{F}(\Psi_2,t)\|\Delta t
  + \|\mathcal{D}_R(\Psi_1)-\mathcal{D}_R(\Psi_2)\|\Delta t \\
  &\le
  \Bigl[
    1 +
    |\Lambda_m|\|\Xi\| L_{\mathcal{F}}(R)\Delta t
    + L_D \Delta t
  \Bigr]\|\Psi_1-\Psi_2\|,
  \end{aligned}
\]
for $\|\Psi_i\|\le R$. Using the multiplicity contraction condition $|\Lambda_m|\|\Xi\|\le \kappa<1$, one can choose $\Delta t$ small enough that
\[
  1 + |\Lambda_m|\|\Xi\| L_{\mathcal{F}}(R)\Delta t + L_D \Delta t
  \le
  1 - \eta
\]
for some $\eta>0$, thereby making $\Phi_{\Delta t}$ a strict contraction on the appropriate bounded subset of $\mathcal{H}$.

\begin{proposition}[Discrete MVF Step as Uniform Contraction]
Under (H1)--(H5) and for sufficiently small $\Delta t>0$, the one-step MVF map $\Phi_{\Delta t}$ is a contraction on bounded subsets of $\mathcal{H}$, i.e.\ there exists $\eta\in(0,1)$ such that
\[
  \|\Phi_{\Delta t}(\Psi_1)-\Phi_{\Delta t}(\Psi_2)\|
  \le
  (1-\eta)\,\|\Psi_1-\Psi_2\|
\]
whenever $\|\Psi_i\|\le R$, with $\eta$ depending on $R$ and $\Delta t$.
\end{proposition}

\begin{proof}
As above, choose $\Delta t$ such that
\[
  |\Lambda_m|\|\Xi\| L_{\mathcal{F}}(R)\Delta t + L_D \Delta t \le -\eta
\]
for some $\eta\in(0,1)$; this is possible provided the coefficients are bounded. Then the desired inequality follows immediately from the previous bound.
\end{proof}

When composed with the projectors $P_E$ and $\Pi_{\mathrm{CSL}}$, which are non-expansive (orthogonal projectors are $1$-Lipschitz), the full step
\[
  \Psi_{t+\Delta t}
  = P_E \Pi_{\mathrm{CSL}}\Phi_{\Delta t}(\Psi_t)
\]
remains a contraction with the same Lipschitz constant. Hence under repeated application, the discrete MVF-augmented MultiplicityCell step converges to a unique fixed point in the lawful-ethical subspace, consistent with the continuous-time convergence analysis.

\section{ACE Budget and Norm Guards}

To ensure compatibility with the PIRTM ACE budget, we make the following explicit.

\begin{description}
  \item[(H7) ACE Budget.] There exists $\varepsilon>0$ such that
        \[
          \sup_t\|\Xi(t)\|\le 1-\varepsilon,
          \quad
          |\Lambda_m(t)|L_T \le c < \varepsilon,
        \]
        where $L_T$ is a Lipschitz bound on the residual map $T_{\Lambda_m}$. This implies
        \[
          |\Lambda_m(t)|\,\|T_{\Lambda_m}\| \le c < \varepsilon,
        \]
        and hence overall one-step contraction $\le 1-\varepsilon+c < 1$.
\end{description}

In the continuous MVF setting, we ensure consistency by requiring that for any step,
\[
  \|\Lambda_m \Xi \mathcal{F}(\Psi,t)\|\Delta t
  \le \kappa (1-\varepsilon)\|\Psi\|
\]
and that the dissipative term $\mathcal{D}_R$ yields additional contraction.

The implementation-level guard
\[
  \|\text{residual}\|\,|\Lambda_m| \le \kappa (1-\varepsilon)
\]
enforces this ACE condition numerically, preventing the emergence of steps that would violate the underlying contraction assumptions and ensuring that MVF operates strictly within the admissible contraction band.

\bigskip

\noindent
\textbf{Summary.} The estimates and proofs in this appendix show that under the boundedness and Lipschitz assumptions used in the main article, MVF:
\begin{itemize}
  \item admits a unique global solution (well-posedness),
  \item has a Lyapunov functional that is nonincreasing and, under mild convexity assumptions, leads to convergence to a unique resonant fixed point, and
  \item can be embedded into the discrete MultiplicityCell recursion via a Lie--Trotter-like split step that is itself a contraction under ACE-calibrated step sizes.
\end{itemize}
These details complete the mathematical underpinning of the claims made in the main text.

\nocite{*}
\bibliographystyle{unsrt}  
\bibliography{references}  


\end{document}
